% ============================================================================
% COURS 13 — MACHINES ASYNCHRONES ET SYNCHRONES
% Partie B : machine asynchrone
% Source : 17-MAS_MS.docx
% ============================================================================

\section{Machine asynchrone}
\label{sec:masms-4}

\subsection{Principe et glissement}

Dans une machine asynchrone, le champ tournant statorique balaie les conducteurs du rotor et y induit des forces électromotrices. Les circuits rotoriques étant fermés, des courants induits apparaissent. L'action du champ statorique sur ces courants crée le couple électromagnétique.

Le rotor ne peut pas atteindre la vitesse du champ : au synchronisme, il ne verrait plus de variation de flux, la f.é.m. rotorique et le couple deviendraient nuls.

On définit le \textbf{glissement} :
\[
\boxed{g=\frac{\Omega_s-\Omega}{\Omega_s}=\frac{N_s-N}{N_s}.}
\]

La fréquence des courants rotoriques est :
\[
\boxed{f_r=g f,}
\qquad
\boxed{\omega_r=g\omega.}
\]

Les principaux domaines de fonctionnement sont :
\begin{center}
\small
\begin{tabularx}{\linewidth}{>{\bfseries}c c X}
\toprule
Glissement & Vitesse & Fonctionnement \\
\midrule
$0<g<1$ & $0<\Omega<\Omega_s$ & moteur.\\
$g<0$ & $\Omega>\Omega_s$ & génératrice ou freinage hypersynchrone.\\
$g>1$ & $\Omega<0$ avec champ direct & freinage par contre-courant.\\
$g=1$ & $\Omega=0$ & démarrage, rotor bloqué.\\
$g=0$ & $\Omega=\Omega_s$ & couple nul.\\
\bottomrule
\end{tabularx}
\end{center}

\begin{TLMachineWarning}{Pertes rotoriques}
Plus le glissement est élevé, plus les pertes Joule rotoriques sont importantes. Le point nominal d'une MAS industrielle se situe généralement à faible glissement.
\end{TLMachineWarning}

\subsection{Schéma monophasé équivalent}

En régime permanent triphasé équilibré, la MAS peut être étudiée par un schéma équivalent monophasé. Les grandeurs rotoriques sont ramenées au stator.

\begin{center}\TLSchemaEquivalentMAS\end{center}

Les éléments du modèle sont :
\begin{itemize}
  \item $R_1$ : résistance d'une phase statorique ;
  \item $X_1=L_1\omega$ : réactance de fuite statorique ;
  \item $R_f$ : résistance modélisant les pertes fer ;
  \item $X_0=L_0\omega$ : réactance de magnétisation ;
  \item $R$ et $X=L\omega$ : résistance et réactance de fuite rotoriques ramenées au stator ;
  \item $R/g$ : résistance fictive traduisant la puissance transmise au rotor.
\end{itemize}

La puissance transmise du stator au rotor vaut :
\[
\boxed{P_{tr}=3\frac{R}{g}{I_1'}^2.}
\]
Les pertes fer statoriques sont modélisées par :
\[
\boxed{P_{fs}=3\frac{V_1^2}{R_f}.}
\]

\TLMAImage[.54\linewidth]{image52.png}

\subsection{Bilan de puissance et rendement}

\begin{center}\TLBilanPuissanceMAS\end{center}

Le bilan de puissance s'écrit :
\[
P_1-P_{Js}-P_{fs}=P_{tr},
\]
\[
\boxed{P_{Jr}=gP_{tr},}
\qquad
\boxed{P_{em}=(1-g)P_{tr}=C_{em}\Omega.}
\]
Après retrait des pertes mécaniques :
\[
P_u=P_{em}-P_{\mathrm{mec}},
\qquad
\boxed{\eta=\frac{P_u}{P_1}.}
\]

La relation $P_{tr}=C_{em}\Omega_s$ est particulièrement utile :
\[
\boxed{C_{em}=\frac{P_{tr}}{\Omega_s}.}
\]

\TLMAImage[.82\linewidth]{image53.png}

\begin{TLMachineExercise}{Bilan nominal d'une MAS}
Une machine porte les indications $230/400\ \mathrm V$, $50\ \mathrm{Hz}$, $15\ \mathrm{kW}$, $1440\ \mathrm{tr\,min^{-1}}$. À charge nominale, le courant de ligne vaut $19{,}4\ \mathrm A$ et la puissance absorbée $11\ \mathrm{kW}$. Les pertes fer et mécaniques sont négligées.
\begin{enumerate}
  \item Déterminer le nombre de paires de pôles et le glissement.
  \item Calculer la puissance réactive et le facteur de puissance.
  \item Déterminer le couple nominal et les pertes Joule rotoriques.
  \item Identifier les paramètres du schéma équivalent à partir d'un essai à vide et d'un essai rotor bloqué.
\end{enumerate}
\end{TLMachineExercise}

\subsection{Couple électromagnétique}

Pour l'étude du couple, les chutes de tension dans $R_1$ et $X_1$ sont souvent négligées. La branche rotorique est alors alimentée sous $V_1$ :
\[
\underline Z_1=\frac{R}{g}+jX,
\qquad
I_1'=\frac{V_1}{\sqrt{(R/g)^2+X^2}}.
\]

En utilisant $P_{tr}=3(R/g){I_1'}^2=C_{em}\Omega_s$ :
\[
\boxed{C_{em}=\frac{3V_1^2}{\Omega_s}\frac{Rg}{R^2+(Xg)^2}.}
\]
Une écriture équivalente est :
\[
C_{em}=\frac{p}{\omega}\,
\frac{3V_1^2(R/g)}{(R/g)^2+X^2}.
\]

\begin{center}\TLCoupleGlissementMAS\end{center}
\TLMAImage[.68\linewidth]{image56.png}

Le couple maximal est obtenu pour :
\[
\boxed{g_M=\frac{R}{X}.}
\]
Sa valeur est indépendante de $R$ :
\[
\boxed{C_M=\frac{3V_1^2}{2\Omega_sX}.}
\]
Au démarrage, $g=1$ :
\[
\boxed{C_D=\frac{3V_1^2}{\Omega_s}\frac{R}{R^2+X^2}.}
\]

Au voisinage du synchronisme, $g\ll g_M$, la caractéristique est presque linéaire :
\[
\boxed{C_{em}\simeq\frac{3V_1^2}{\Omega_sR}\,g.}
\]

\begin{TLMachineMethod}{Étudier une MAS à partir du schéma équivalent}
\begin{enumerate}
  \item Déterminer $p$, $\Omega_s$ et $g$.
  \item Identifier les paramètres par essai à vide et rotor bloqué.
  \item Calculer $P_{tr}$ puis $C_{em}=P_{tr}/\Omega_s$.
  \item Rechercher $g_M$ et $C_M$.
  \item Placer le point de fonctionnement par l'intersection entre couple moteur et couple résistant.
  \item Vérifier les courants, le rendement et la stabilité.
\end{enumerate}
\end{TLMachineMethod}

\begin{TLMachineExercise}{Expression et maximum du couple}
À partir du schéma simplifié comportant $R/g$ et $l\omega$, montrer que :
\[
C=\frac{3pV^2}{\omega}\frac{R/g}{(R/g)^2+(l\omega)^2}.
\]
Déterminer $g_M$, $C_M$, la vitesse correspondante et tracer la caractéristique pour une machine dont la vitesse de synchronisme vaut $1500\ \mathrm{tr\,min^{-1}}$.
\end{TLMachineExercise}

\subsection{Variation de vitesse}

La vitesse mécanique est :
\[
\boxed{\Omega=\frac{\omega}{p}(1-g).}
\]
Elle peut être modifiée en agissant sur :
\begin{itemize}
  \item le glissement, notamment par la tension ;
  \item la fréquence statorique ;
  \item le nombre de paires de pôles.
\end{itemize}

La solution moderne consiste à alimenter la machine par un onduleur et à modifier la fréquence. Pour conserver le flux magnétique approximativement constant, on maintient :
\[
\boxed{\frac{V_1}{f}=\mathrm{constante}.}
\]

\begin{center}\TLCommandeVsurF\end{center}

Sous la fréquence nominale, le flux et le couple maximal restent presque constants. Au-delà de la fréquence nominale, la tension ne peut plus augmenter : le fonctionnement se poursuit à puissance approximativement constante et le couple disponible diminue.

\TLMAImage[.56\linewidth]{image58.png}
\TLMAImage[.56\linewidth]{image59.png}

Dans la zone utile et pour un faible glissement, la caractéristique peut s'écrire sous forme affine :
\[
C_{em}\simeq K(\Omega_s-\Omega).
\]
Cette propriété explique la translation des courbes couple--vitesse lorsque la fréquence varie.

\begin{TLMachineWarning}{Limites de la commande scalaire}
La loi $V/f=\mathrm{cte}$ compense correctement le flux à vitesse moyenne et élevée. À basse fréquence, la chute de tension statorique devient comparable à la tension appliquée : une compensation ou une commande vectorielle est alors nécessaire.
\end{TLMachineWarning}

\begin{TLMachineExercise}{Identification et commande $V/f$}
Une MAS $230/400\ \mathrm V$, $50\ \mathrm{Hz}$, $1440\ \mathrm{tr\,min^{-1}}$ est caractérisée par :
\begin{itemize}
  \item essai à vide : $V=230\ \mathrm V$, $I_0=9{,}5\ \mathrm A$, $P_0=630\ \mathrm W$ ;
  \item essai rotor bloqué : $V_{cc}=50\ \mathrm V$, $I_{cc}=30\ \mathrm A$, $P_{cc}=830\ \mathrm W$.
\end{itemize}
Déterminer les paramètres du modèle, établir l'expression du couple et montrer l'effet d'une commande $V/f$ constante sur les caractéristiques couple--vitesse.
\end{TLMachineExercise}

\subsection{Fonctionnement dans les quatre quadrants}

Une association onduleur--MAS réversible peut fonctionner en moteur ou en génératrice dans les deux sens de rotation. Le signe de la puissance mécanique est donné par :
\[
P_m=C\Omega.
\]

\begin{center}\TLQuadrantsMachine\end{center}

La permutation de deux phases inverse le sens du champ tournant. Pour permettre le freinage régénératif, la chaîne d'énergie doit être réversible jusqu'à la source ; sinon, l'énergie est dissipée dans une résistance de freinage.
